% =====================================================================
%  c00.tex —— 封面、译者说明、目录（Lead 独占）
% =====================================================================
\begin{titlepage}
\centering
\vspace*{2.2cm}
{\Large 学习讨论班}\\[1.0em]
{\Huge\bfseries SYZ 镜像对称}\\[0.7em]
{\large Learning Seminar: SYZ Mirror Symmetry}\\[2.2em]
{\Large 中文译注本}\\[1.6em]
{\normalsize
讨论班报告人：Si--Yang Liu、Yunpeng Niu、Kyoungmo Kim、Maximilian Schimpf\\[0.6em]
底本：\texttt{https://diewildkatze.com/SYZ.html}\\[0.6em]
译本：2026 年 11 月}
\vfill
{\small 本译本仅供学习交流；原文版权归原作者所有。\\
公式与定理均按底本 HTML 中的 \LaTeX{} 源码逐字转写，未改动数学内容；\\
译者的补充说明一律写在 \texttt{\% [Tutor Note: ...]} 注释中，不进入排版。}
\end{titlepage}

\thispagestyle{empty}
\section*{译者说明}
\addcontentsline{toc}{section}{译者说明}

本译本由底本网页（一份关于 SYZ 镜像对称的学习讨论班讲义）逐段译出，工程规范见交付目录中的
\texttt{README.md}。几点需要先说明：

\begin{enumerate}[leftmargin=1.6em,itemsep=2pt]
  \item \textbf{底本形态}：底本不是 PDF，而是网页（HTML + MathJax）。因此不存在「扫描页图像」，
        公式的\textbf{唯一真值}是网页 HTML 中的 \LaTeX{} 源码；这是本译本相比一般 PDF 翻译的优势
        （无需从图像猜符号），但也意味着底本自身的笔误会被原样带入，故设「原文笔误登记表」。
  \item \textbf{公式编号}：底本没有任何编号公式，本译本同样\textbf{不编号}（一律
        \verb|\[ ... \]| 或 \verb|equation*|）；正文中的引用一律为纯文本。
  \item \textbf{图表}：底本无可下载的插图（\texttt{<img>} 标签数为 0），有 3 张表、3 个由
        \texttt{q.uiver.app} 嵌入的交换图（本译本用 \texttt{tikz-cd} 重画）。
  \item \textbf{参考文献}：按学术惯例保留英文原貌，仅转录，不逐条校对（见 README「已知限制」）。
  \item \textbf{批注}：以 \texttt{\% [Tutor Note: ...]} 形式写在源码里，编译后不显示；
        全部批注汇总在交付目录的「译注汇总.md」。
\end{enumerate}

\clearpage
\tableofcontents
\clearpage
